\setpart{Pi0.5 · 测量结果}

\begin{frame}[t]{同协议下的端到端延迟与单连接可持续周期}
\vspace{-.1cm}
\small\tb{单请求端到端延迟，p50（毫秒），RTX 5880，OpenPI WebSocket，10 步}\\[.08cm]
\begin{tikzpicture}[x=1cm,y=.5cm]
  \def\k{.034}
  \foreach \y/\name/\mid/\colr in {4/M* 默认前端（Rust）/66/Teal, 3/SGLang core（分桶）/97/Blue, 2/SGLang core（默认）/167/Blue, 1/vLLM-Omni/211/Orange}{
    \node[anchor=east,font=\scriptsize,inner sep=1pt] at (2.3,\y) {\name};
    \fill[\colr!70] (2.4,\y-.22) rectangle ({2.4+\mid*\k},\y+.22);
    \node[anchor=west,font=\scriptsize\bfseries,inner sep=1pt] at ({2.42+\mid*\k},\y) {\mid};
  }
\end{tikzpicture}\\[.05cm]
\scriptsize
\renewcommand{\arraystretch}{1.1}
\begin{tabularx}{\textwidth}{>{\bfseries}p{3.4cm} Y Y Y}
\toprule
\rowcolor{Mist}
 & 单连接可持续的最短周期 & 周期 100 毫秒时的表现 & 备注 \\
\midrule
M*（Rust 前端） & 50--75 毫秒 & 每秒 10.0 个请求，无跳过 & 最新版 2c8ffde8 \\
SGLang core（分桶 / 默认） & 约 100 毫秒 / 150--200 毫秒 & 分桶：跳过约 5\% 的触发；默认：排队 & 默认的图按精确长度缓存，只有 4 个名额 \\
vLLM-Omni & 200--300 毫秒 & 每秒 3.1 个请求，跳过 275 个触发 & eager 模式 \\
\bottomrule
\end{tabularx}
\par\takeaway{同一协议下，SGLang 的默认配置不是它的最佳配置。比较时取分桶后的数字，M* 与 vLLM-Omni 的差距在 3 倍左右。}
\par\memo{来源：设计文档 §2.6.1（138 个请求）、§2.6.2（每档 40 个请求）、§2.6.3；runs/2026-11-20\_\allowbreak{}ws-probe/\allowbreak{}；runs/2026-11-21\_\allowbreak{}sglang-buckets-base/\allowbreak{}。}
\end{frame}

\begin{frame}[t]{请求断开不会中止在途推理：新请求等的是剩余时间}
\vspace{-.1cm}
\begin{tikzpicture}[x=1cm,y=.45cm]
  \def\X{2.4}
  \def\S{.075}
  % 时间刻度（毫秒）
  \foreach \tm in {0,50,100,150}{
    \draw[Line,line width=.4pt] ({\X+\tm*\S},-2.5) -- ({\X+\tm*\S},.9);
    \node[font=\tiny,text=Muted] at ({\X+\tm*\S},-2.8) {\tm\,毫秒};
  }
  \node[anchor=west,font=\scriptsize\bfseries,text=Ink] at (0,.5) {请求 A};
  
  \fill[Teal!55] ({\X},.05) rectangle ({\X+66.2*\S},.5);
  \node[font=\tiny,text=white,anchor=center] at ({\X+33*\S},.28) {A：66 毫秒};
  \node[anchor=west,font=\scriptsize\bfseries,text=Ink] at (0,-1.1) {请求 B};
  
  \draw[Red,densely dashed,line width=.7pt] ({\X+26.5*\S},.9) -- ({\X+26.5*\S},-1.6);
  \node[font=\tiny,text=Red,anchor=south west] at ({\X+26.5*\S+.02},-1.45) {B 到达，26.5 毫秒};
  \fill[Orange!55] ({\X+66.2*\S},-1.3) rectangle ({\X+93.2*\S},-1.75);
  \fill[Muted!30] ({\X+26.5*\S},-1.3) rectangle ({\X+66.2*\S},-1.75);
  \node[font=\tiny,text=Ink,anchor=center] at ({\X+46*\S},-1.52) {等 A 的剩余推理};
  \node[font=\tiny,text=Orange,anchor=west] at ({\X+93.2*\S+.05},-1.52) {B 完成：93 毫秒};
  \draw[Blue,densely dashed,line width=.7pt] ({\X+66.2*\S},.9) -- ({\X+66.2*\S},-2.5);
\end{tikzpicture}\\[.05cm]
\scriptsize
\renewcommand{\arraystretch}{1.1}
\begin{tabularx}{\textwidth}{>{\bfseries}p{2.6cm} Y Y Y Y}
\toprule
\rowcolor{Mist}
服务（毫秒，p50） & 对照（单请求） & B，A 关闭连接 & B，A 保持打开 & 预测（不中止） \\
\midrule
M*（Rust 前端） & 66.2 & 93.2 & 93.2 & 105.9 \\
SGLang core & 164.7 & 258.5 & 257.1 & 263.5 \\
vLLM-Omni & 210.4 & 336.4 & 328.2 & 336.7 \\
\bottomrule
\end{tabularx}
\par\takeaway{断开后，服务端仍把 A 算完；如果会中止，B 应接近对照值。测得值与“不中止”的预测相符（SGLang、vLLM 偏差在 3\% 内；M* 低 9--12\%，原因未查明）。}
\par\memo{来源：设计文档 §2.6.4（修正版探测，每框架 15 组，runs/2026-11-20\_\allowbreak{}ws-probe-cancel-v2/\allowbreak{}）。第一版探测曾把 Rust 前端误判为会中止，已修正。}
\end{frame}

\begin{frame}[t]{Level 1：各触发策略下的欠载占比（多频输入，对齐执行）}
\vspace{-.1cm}
\scriptsize
\renewcommand{\arraystretch}{1.15}
\setlength{\tabcolsep}{4pt}
\begin{tabular}{>{\bfseries}p{3.4cm} c c c c}
\toprule
\rowcolor{Mist}
触发策略 & \textcolor{Teal}{M*} & \textcolor{Orange}{vLLM-Omni} & \textcolor{Blue}{SGLang（默认）} & \textcolor{Blue}{SGLang（分桶）} \\
\midrule
周期 100 毫秒 & \colorbox{Mist}{\makebox[1.7cm][c]{0.0\%}} & \colorbox{RedBg}{\makebox[1.7cm][c]{0.5\%}} & \colorbox{Mist}{\makebox[1.7cm][c]{0.0\%}} & \colorbox{Mist}{\makebox[1.7cm][c]{0.0\%}} \\
周期 200 毫秒 & \colorbox{Mist}{\makebox[1.7cm][c]{0.0\%}} & \colorbox{RedBg}{\makebox[1.7cm][c]{3.5\%}} & \colorbox{Mist}{\makebox[1.7cm][c]{0.0\%}} & \colorbox{Mist}{\makebox[1.7cm][c]{0.0\%}} \\
周期 300 毫秒 & \colorbox{Mist}{\makebox[1.7cm][c]{0.0\%}} & \colorbox{RedBg}{\makebox[1.7cm][c]{2.2\%}} & \colorbox{Mist}{\makebox[1.7cm][c]{0.0\%}} & \colorbox{Mist}{\makebox[1.7cm][c]{0.0\%}} \\
每 8 步重规划 & \colorbox{Mist}{\makebox[1.7cm][c]{0.0\%}} & \colorbox{Red!35}{\makebox[1.7cm][c]{21.4\%}} & \colorbox{Mist}{\makebox[1.7cm][c]{0.0\%}} & \colorbox{Mist}{\makebox[1.7cm][c]{0.0\%}} \\
阈值 0.3 秒 & \colorbox{Mist}{\makebox[1.7cm][c]{0.0\%}} & \colorbox{RedBg}{\makebox[1.7cm][c]{4.8\%}} & \colorbox{Mist}{\makebox[1.7cm][c]{0.0\%}} & \colorbox{Mist}{\makebox[1.7cm][c]{0.0\%}} \\
阈值 0.5 秒 & \colorbox{Mist}{\makebox[1.7cm][c]{0.0\%}} & \colorbox{RedBg}{\makebox[1.7cm][c]{2.3\%}} & \colorbox{Mist}{\makebox[1.7cm][c]{0.0\%}} & \colorbox{Mist}{\makebox[1.7cm][c]{0.0\%}} \\
阈值 0.8 秒 & \colorbox{Mist}{\makebox[1.7cm][c]{0.0\%}} & \colorbox{RedBg}{\makebox[1.7cm][c]{5.0\%}} & \colorbox{Mist}{\makebox[1.7cm][c]{0.0\%}} & \colorbox{Mist}{\makebox[1.7cm][c]{0.0\%}} \\
\bottomrule
\end{tabular}
\par\takeaway{M* 与分桶 SGLang 在所有策略下都不欠载。vLLM-Omni 的欠载来自慢回复：块只有 1 秒，一次慢回复就用掉大部分缓冲。}
\par\memo{来源：runs/2026-11-20\_\allowbreak{}level1/\allowbreak{}level1\_\allowbreak{}summary.md（三轮中位数，欠载 tick 占比）；设计文档 §2.7。控制器与输入都是模拟的。}
\end{frame}

\begin{frame}[t]{动作有多新：动作年龄随端到端延迟上升}
\vspace{-.1cm}
\small\tb{周期 200 毫秒，多频输入，对齐执行（毫秒，p50）}\\[.08cm]
\begin{tikzpicture}[x=1cm,y=.4cm]
  \def\k{.017}
  \foreach \y/\name/\e/\a/\colr in {4/M*（Rust）/61/134/Teal, 3/SGLang（分桶）/95/201/Blue, 2/SGLang（默认）/110/266/Blue, 1/vLLM-Omni/227/469/Orange}{
    \node[anchor=east,font=\scriptsize,inner sep=1pt] at (2.3,\y+.12) {\name};
    \fill[\colr!30] (2.4,\y+.05) rectangle ({2.4+\e*\k},\y+.32);
    \node[anchor=west,font=\tiny,inner sep=1pt] at ({2.42+\e*\k},\y+.18) {端到端 \e};
    \fill[\colr!80] (2.4,\y-.32) rectangle ({2.4+\a*\k},\y-.05);
    \node[anchor=west,font=\tiny\bfseries,inner sep=1pt] at ({2.42+\a*\k},\y-.18) {动作年龄 \a};
  }
\end{tikzpicture}\\[.05cm]
\scriptsize
\begin{nbox}{Gold}{朴素执行（OpenPI 参考客户端的做法）}
\scriptsize 每个动作比计划时刻晚约一个回复延迟：M* 晚 68 毫秒，SGLang（分桶）晚 135 毫秒，vLLM-Omni 晚 337 毫秒。以每 8 步重规划为例：朴素执行的 vLLM-Omni 欠载只有 2.3\%，对齐执行时为 21.4\%。朴素执行是用更旧的动作填住缺口换来的。
\end{nbox}
\par\takeaway{动作年龄由触发间隔和端到端延迟共同决定。要降低动作年龄，只能缩短端到端延迟或加快触发；取消旧请求做不到（见断开实验）。}
\par\memo{来源：runs/2026-11-20\_\allowbreak{}level1/\allowbreak{}level1\_\allowbreak{}summary.md；设计文档 §2.7 的表与第 5 条。}
\end{frame}
